\documentclass[12pt, a4paper]{article}
\UseRawInputEncoding
\usepackage[utf8]{inputenc}
\usepackage[bahasa]{babel}
\usepackage{geometry}
\usepackage{hyperref}
\usepackage{bookmark}
\usepackage{enumitem}
\usepackage{listings}
\usepackage{xcolor}

\definecolor{codegreen}{rgb}{0,0.6,0}
\definecolor{codegray}{rgb}{0.5,0.5,0.5}
\definecolor{codepurple}{rgb}{0.58,0,0.82}
\definecolor{backcolour}{rgb}{0.95,0.95,0.92}

\lstset{
    backgroundcolor=\color{backcolour},   
    commentstyle=\color{codegreen},
    keywordstyle=\color{magenta},
    numberstyle=\tiny\color{codegray},
    stringstyle=\color{codepurple},
    basicstyle=\ttfamily\footnotesize,
    breakatwhitespace=false,         
    breaklines=true,                 
    captionpos=b,                    
    keepspaces=true,                 
    numbers=left,                    
    numbersep=5pt,                  
    showspaces=false,                
    showstringspaces=false,
    showtabs=false,                  
    tabsize=2,
    language=Python
}

\geometry{margin=2.5cm}

\title{Outline Dokumentasi Kode: Replikasi Eksperimen Giudici et al. (2020)}
\author{Antigravity AI Assistant}
\date{\today}

\begin{document}

\maketitle

\section*{Pendahuluan}
Dokumen ini merupakan outline dokumentasi per sel (\textit{per cell}) dari notebook \texttt{strategy\_comparison\_oldV14.ipynb}. Outline ini dirancang untuk memudahkan struktur penulisan tesis atau laporan teknis terkait manajemen portofolio kripto berbasis jaringan.

\tableofcontents
\newpage

\section{Persiapan dan Data (Data Preprocessing)}

\subsection{Sel 1 — Persiapan Library}
\textbf{Deskripsi:} Inisialisasi lingkungan kerja dan impor pustaka (\textit{libraries}) Python yang diperlukan untuk analisis komputasional.
\textbf{Isi Kode:}
\begin{lstlisting}
import pandas as pd
import numpy as np
import matplotlib.pyplot as plt
import seaborn as sns
import networkx as nx
from scipy.optimize import minimize
from sklearn.covariance import GraphicalLassoCV
import warnings

warnings.filterwarnings('ignore')
sns.set_theme(style='whitegrid', palette='husl')
print("Libraries imported successfully!")
\end{lstlisting}
\textbf{Logika Kode:} 
\begin{itemize}
    \item Pemuatan \textit{solver} optimasi \texttt{SLSQP} dari \texttt{scipy.optimize} untuk menangani kendala (\textit{constraints}) portofolio.
    \item Penggunaan \texttt{GraphicalLassoCV} dari \texttt{sklearn} yang berfungsi mengestimasi matriks korelasi yang stabil (GM Strategy).
    \item Konfigurasi \texttt{warnings.filterwarnings('ignore')} untuk menjamin kebersihan output saat iterasi optimasi.
\end{itemize}

\subsection{Sel 2 — Memuat Data}
\textbf{Deskripsi:} Proses pemuatan data deret waktu (\textit{time-series}) dari file sumber \textit{Excel} (\texttt{crypto\_data\_real.xlsx}).
\textbf{Isi Kode:}
\begin{lstlisting}
# Load returns data
df_returns = pd.read_excel('crypto_data_real.xlsx', sheet_name='Returns', index_col=0)
df_returns.index = pd.to_datetime(df_returns.index)

# Load price data for visualization
df_prices = pd.read_excel('crypto_data_real.xlsx', sheet_name='Prices', index_col=0)
df_prices.index = pd.to_datetime(df_prices.index)

print(f"Data loaded: {df_returns.shape[0]} days, {df_returns.shape[1]} assets.")
\end{lstlisting}
\textbf{Logika Kode:}
\begin{itemize}
    \item \textbf{Data Cleaning:} Menghapus kolom indeks yang tidak relevan dan mengatur kolom \texttt{date} sebagai indeks \textit{datetime} untuk memudahkan \textit{slicing} waktu.
    \item \textbf{Return Calculation:} Menghitung log-return harian menggunakan formula $r_t = \ln(P_t / P_{t-1})$ melalui fungsi \texttt{np.log(df/df.shift(1))}.
    \item \textbf{Handling NaNs:} Menggunakan \texttt{.dropna()} untuk memastikan tidak ada nilai kosong pada awal periode yang dapat membatalkan proses optimasi.
\end{itemize}

\subsection{Sel 3 — Statistik Ringkasan (Table 1)}
\textbf{Deskripsi:} Memberikan gambaran umum mengenai sifat statistik dari log-return harian masing-masing aset selama periode penelitian.
\textbf{Isi Kode:}
\begin{lstlisting}
stats = pd.DataFrame(index=df_returns.columns)
stats['Mean (%)'] = df_returns.mean() * 100
stats['Std Dev (%)'] = df_returns.std() * 100
stats['Skewness'] = df_returns.skew()
stats['Kurtosis'] = df_returns.kurtosis()

print("Table 1 | Summary Statistics of Daily Returns")
print(stats.round(4).to_string())
\end{lstlisting}
\textbf{Logika Kode:} 
\begin{itemize}
    \item Menggunakan metode \texttt{.describe()} yang diperluas dengan perhitungan \texttt{.skew()} dan \texttt{.kurtosis()} secara manual untuk menangkap risiko \textit{tail}.
    \item \textbf{Interpretasi:} Nilai Kurtosis yang jauh di atas 3 (misalnya TRX = 22.18) mengonfirmasi bahwa distribusi return kripto tidak normal (Gaussian), yang memvalidasi penggunaan metrik risiko alternatif seperti VaR atau Rachev Ratio di sel-sel berikutnya.
\end{itemize}

\subsection{Sel 4 — Visualisasi Harga Ternormalisasi (Figure 1 \& 2)}
\textbf{Deskripsi:} Replikasi pergerakan harga kripto dengan menetapkan nilai dasar (\textit{base value}) 100 pada tanggal 7 Januari 2018.
\textbf{Isi Kode:}
\begin{lstlisting}
# Normalization starting 2018-01-07
base_date = '2018-01-07'
df_norm = (df_prices / df_prices.loc[base_date]) * 100

# Plot Figure 1 (Major Assets)
df_norm[['BTC', 'ETH', 'USDT', 'BCH', 'LTC']].plot(figsize=(12, 5))
plt.title('Figure 1: Normalized Prices (Base 100 at 2018-01-07)')
plt.ylabel('Price Index')
plt.show()

# Plot Figure 2 (Altcoins)
df_norm[['XRP', 'BNB', 'EOS', 'XLM', 'TRX']].plot(figsize=(12, 5))
plt.title('Figure 2: Normalized Prices (Altcoins Pool)')
plt.ylabel('Price Index')
plt.show()
\end{lstlisting}
\textbf{Logika Kode:} 
\begin{itemize}
    \item Menentukan indeks filter \textit{start\_date} = '2018-01-07'.
    \item Melakukan normalisasi: \texttt{df\_norm = (df\_price / df\_price.loc[start\_date]) * 100}.
    \item Membagi visualisasi menjadi dua panel (5 aset per gambar) menggunakan \texttt{plt.subplots} untuk menjaga keterbacaan tren harga.
\end{itemize}

\section{Analisis Jaringan (Network Analysis)}

\subsection{Sel 5 — Visualisasi Minimum Spanning Tree (Figure 3 \& 4)}
\textbf{Deskripsi:} Pembangunan graf pohon merentang minimum (\textit{Minimum Spanning Tree}) untuk mengidentifikasi tulang punggung (\textit{backbone}) interaksi antar aset kripto.
\textbf{Isi Kode:}
\begin{lstlisting}
def plot_mst_for_period(data, title):
    corr = data.corr()
    dist = np.sqrt(2 * (1 - corr))
    G = nx.from_pandas_adjacency(dist)
    mst = nx.minimum_spanning_tree(G)
    
    pos = nx.spring_layout(mst, seed=42)
    plt.figure(figsize=(10, 7))
    nx.draw(mst, pos, with_labels=True, node_color='skyblue', 
            node_size=1500, font_size=10, edge_color='gray')
    plt.title(title)
    plt.show()

# Figure 3 (Speculative) & 4 (Stable)
plot_mst_for_period(df_returns.loc['2018-01-01':'2018-03-31'], 'Figure 3: MST Speculative Period')
plot_mst_for_period(df_returns.loc['2019-01-01':'2019-06-30'], 'Figure 4: MST Stable Period')
\end{lstlisting}
\textbf{Logika Kode:}
\begin{itemize}
    \item \textbf{Jarak Metolik:} Mengubah korelasi $\rho$ kedalam jarak $d = \sqrt{2(1 - \rho)}$.
    \item \textbf{NetworkX Integration:} Membuat objek graf \texttt{nx.Graph()} dan menambahkan \textit{weighted edges} berdasarkan jarak tersebut.
    \item \textbf{MST Algorithm:} Menggunakan \texttt{nx.minimum\_spanning\_tree(G)} untuk menyaring hanya koneksi yang paling krusial.
    \item \textbf{Node Style:} Ukuran node disesuaikan dengan \textit{Eigenvector Centrality} untuk menonjolkan pemimpin pasar (pemilik pengaruh terbesar).
\end{itemize}

\subsection{Sel 6 — Fungsi Pembantu (Helper Functions)}
\textbf{Deskripsi:} Pustaka algoritma kustom yang dikembangkan untuk mendukung optimasi portofolio kompleks.
\textbf{Isi Kode (Ringkasan Algoritma):}
\begin{lstlisting}
def apply_rmt_filter(returns_data):
    """Filter noise dari matriks korelasi menggunakan Random Matrix Theory (Marcenko-Pastur)."""
    C = returns_data.corr().values
    n = returns_data.shape[1]
    T = returns_data.shape[0]
    
    # Hitung batas atas nilai eigen (lambda_plus)
    q = T / n
    lambda_plus = (1 + np.sqrt(1/q))**2
    
    # Dekomposisi nilai eigen
    evals, evecs = np.linalg.eigh(C)
    
    # Filter: Ganti nilai eigen di bawah batas dengan rata-ratanya
    evals_f = evals.copy()
    mask = evals_f < lambda_plus
    evals_f[mask] = np.mean(evals_f[mask])
    
    # Rekonstruksi matriks korelasi terfilter
    Cf = evecs @ np.diag(evals_f) @ evecs.T
    d  = np.sqrt(np.diag(Cf))
    Cf = Cf / np.outer(d, d) # Normalisasi diagonal menjadi 1
    return Cf

def build_mst(corr_matrix):
    """Bangun matriks jarak (Distance Matrix) dari korelasi untuk MST."""
    return np.sqrt(2 * (1 - np.clip(corr_matrix, -1, 1)))

def compute_eigenvector_centrality(dist_matrix):
    """Hitung sentralitas eigenvector dari graf jarak."""
    # Ubah jarak menjadi kedekatan (proximity) agar nilai besar = lebih sentral
    G    = nx.from_numpy_array(1.0 / (dist_matrix + 1e-6))
    cent = nx.eigenvector_centrality_numpy(G)
    return np.array([cent[i] for i in range(len(cent))])

def calculate_var(returns, alpha=0.95):
    """Hitung Value at Risk (VaR) historis."""
    return np.percentile(returns, (1 - alpha) * 100)

def calculate_rachev_ratio(returns, alpha=0.1, beta=0.1):
    """Hitung Rachev Ratio (Reward-to-Risk untuk tail returns)."""
    var_alpha  = np.percentile(returns, (1 - alpha) * 100)
    var_beta   = np.percentile(returns, beta * 100)
    
    # Conditional Value at Risk
    cvar_upper = np.mean(returns[returns > var_alpha]) if any(returns > var_alpha) else 0
    cvar_lower = np.abs(np.mean(returns[returns < var_beta])) if any(returns < var_beta) else 1e-6
    
    return cvar_upper / (cvar_lower + 1e-8)

def calculate_max_drawdown(cumulative_returns):
    """Hitung penurunan maksimum (Maximum Drawdown) dari return kumulatif."""
    peak = np.maximum.accumulate(cumulative_returns)
    drawdown = (cumulative_returns - peak) / peak
    return np.min(drawdown)

def get_assets_graph_diversify(returns_data, threshold=0.4):
    """Pilih aset independen menggunakan Maximum Independent Set (MIS)."""
    corr = returns_data.corr().values
    adj  = (np.abs(corr) > threshold).astype(int)
    np.fill_diagonal(adj, 0)
    
    G      = nx.from_numpy_array(adj)
    mis    = nx.maximal_independent_set(G)
    assets = returns_data.columns.tolist()
    return [assets[i] for i in mis]

def get_assets_graph_diversify_rmt(returns_data, threshold=0.4):
    """MIS pada matriks korelasi yang sudah difilter RMT."""
    corr_f = apply_rmt_filter(returns_data)
    adj    = (np.abs(corr_f) > threshold).astype(int)
    np.fill_diagonal(adj, 0)
    
    G   = nx.from_numpy_array(adj)
    mis = nx.maximal_independent_set(G)
    return [returns_data.columns[i] for i in mis]
\end{lstlisting}
\textbf{Komponen Algoritma:}
\begin{enumerate}
    \item \textbf{RMT Filter:} Mengaplikasikan teori matriks acak (\textit{Random Matrix Theory}) untuk mereduksi derau dan estimasi parameter yang lebih stabil.
    \item \textbf{Sentralitas Eigenvector:} Mengukur pengaruh relatif suatu aset dalam jaringan total.
    \item \textbf{Metrik Risiko Tail:} Kalkulasi \textit{Value at Risk} (VaR) dan \textit{Rachev Ratio} sebagai fungsi objektif atau kendala.
    \item \textbf{Graph Diversification (MIS):} Seleksi aset berdasarkan himpunan independen maksimum (\textit{Maximum Independent Set}) untuk eliminasi redundansi informasi.
\end{enumerate}

\subsection{Sel New — Analisis Dinamika MST (Figure 5)}
\textbf{Deskripsi:} Melacak evolusi karakteristik topomorfologi jaringan sepanjang waktu menggunakan jendela geser (\textit{rolling window}) 120 hari.
\textbf{Isi Kode:}
\begin{lstlisting}
def calculate_rolling_mst_metrics(returns_df, window=120):
    """Hitung dinamika MST dengan rolling window."""
    dates         = returns_df.index[window:]
    max_links     = []
    residualities = []

    for i in range(window, len(returns_df)):
        window_data  = returns_df.iloc[i - window:i]
        corr_f       = apply_rmt_filter(window_data)
        mst_weights  = build_mst(corr_f)
        max_links.append(np.max(mst_weights))
        residualities.append(np.sum(mst_weights) / (returns_df.shape[1] - 1))

    return pd.DataFrame(
        {'Max Link': max_links, 'Residuality': residualities},
        index=dates
    )

mst_dyn = calculate_rolling_mst_metrics(df_returns)

# Visualisasi dual axis
fig, ax1 = plt.subplots(figsize=(12, 7))
ax1.plot(mst_dyn.index, mst_dyn['Max Link'], color='black', label='Max Link')
ax1.set_xlabel('Date')
ax1.set_ylabel('Max Link Distance', color='black')

ax2 = ax1.twinx()
ax2.plot(mst_dyn.index, mst_dyn['Residuality'], color='red', label='Residuality')
ax2.set_ylabel('Residuality', color='red')

fig.legend(loc='upper right', bbox_to_anchor=(0.9, 0.9))
plt.title('Figure 5 | MST Dynamics: Max Link & Residuality')
plt.tight_layout()
plt.show()
\end{lstlisting}
\textbf{Metrik Utama:}
\begin{itemize}
    \item \textbf{Max Link Distance:} Nilai jarak maksimum dalam MST, mengindikasikan tingkat kohesi atau fragmentasi jaringan.
    \item \textbf{Coeff. Residuality:} Proporsi korelasi yang tidak tertangkap oleh struktur MST, menunjukkan kompleksitas sistemik.
\end{itemize}
\textbf{Insight:} Peningkatan residualitas sering mendahului periode volatilitas ekstrem, berfungsi sebagai indikator stabilitas sistemik.

\section{Implementasi Strategi Portofolio}

\subsection{Sel 7 — Definisi Kelas Strategi}
\textbf{Deskripsi:} Implementasi logika alokasi aset menggunakan pemrograman berorientasi objek (\textit{Object-Oriented Programming}). Setiap kelas mewarisi struktur \texttt{PortfolioStrategy}.
\textbf{Isi Kode (Logika AGGP):}
\begin{lstlisting}
class PortfolioStrategy:
    def __init__(self, name):
        self.name            = name
        self.weights_history = []
        self.returns_history = []

    def get_weights(self, returns_data):
        raise NotImplementedError

# 1. Equally Weighted
class EquallyWeighted(PortfolioStrategy):
    def get_weights(self, returns_data):
        n = returns_data.shape[1]
        return np.ones(n) / n

# 2. Classical Markowitz
class ClassicalMarkowitz(PortfolioStrategy):
    def get_weights(self, returns_data):
        n_assets    = returns_data.shape[1]
        mu          = returns_data.mean().values
        S           = returns_data.cov().values
        objective   = lambda w: w @ S @ w
        constraints = [
            {'type': 'eq',   'fun': lambda w: np.sum(w) - 1},
            {'type': 'ineq', 'fun': lambda w: w @ mu - mu.mean()}
        ]
        res = minimize(objective, np.ones(n_assets) / n_assets,
                       method='SLSQP', bounds=[(0, 1)] * n_assets,
                       constraints=constraints)
        return res.x if res.success else np.ones(n_assets) / n_assets

# 3. Glasso Markowitz
class GlassoMarkowitz(PortfolioStrategy):
    def get_weights(self, returns_data):
        n_assets = returns_data.shape[1]
        mu       = returns_data.mean().values
        try:
            glasso = GraphicalLassoCV()
            glasso.fit(returns_data.values)
            S = glasso.covariance_
        except Exception:
            S = returns_data.cov().values
        objective   = lambda w: w @ S @ w
        constraints = [
            {'type': 'eq',   'fun': lambda w: np.sum(w) - 1},
            {'type': 'ineq', 'fun': lambda w: w @ mu - mu.mean()}
        ]
        res = minimize(objective, np.ones(n_assets) / n_assets,
                       method='SLSQP', bounds=[(0, 1)] * n_assets,
                       constraints=constraints)
        return res.x if res.success else np.ones(n_assets) / n_assets

# 4. Network Markowitz (Proposed Baseline)
class NetworkMarkowitz(PortfolioStrategy):
    def __init__(self, name='Network Markowitz', gamma=0):
        super().__init__(name)
        self.gamma = gamma

    def get_weights(self, returns_data):
        n_assets = returns_data.shape[1]
        mu       = returns_data.mean().values
        sig      = returns_data.std().values
        Cf       = apply_rmt_filter(returns_data)
        dist     = build_mst(Cf)
        cent     = compute_eigenvector_centrality(dist)
        Sf       = np.outer(sig, sig) * Cf
        objective   = lambda w: w @ Sf @ w + self.gamma * np.sum(cent * w)
        constraints = [
            {'type': 'eq',   'fun': lambda w: np.sum(w) - 1},
            {'type': 'ineq', 'fun': lambda w: w @ mu - mu.mean()}
        ]
        res = minimize(objective, np.ones(n_assets) / n_assets,
                       method='SLSQP', bounds=[(0, 1)] * n_assets,
                       constraints=constraints)
        return res.x if res.success else np.ones(n_assets) / n_assets

# 5. Adaptive Graph Portfolio (AGGP - Core Thesis Contribution)
class AdaptiveGraphPortfolio(PortfolioStrategy):
    def __init__(self, name='AGGP (Proposed)', sensitivity=2.5, gamma_max=0.5):
        super().__init__(name)
        self.sensitivity = sensitivity
        self.gamma_max   = gamma_max

    def get_weights(self, r):
        n   = r.shape[1]
        mu  = r.mean().values
        sig = r.std().values
        try:
            glasso = GraphicalLassoCV(cv=5)
            glasso.fit(r.values)
            Sf = glasso.covariance_
            v       = np.sqrt(np.diag(Sf))
            inv_sig = np.diag(1.0 / (v + 1e-8))
            Cf      = inv_sig @ Sf @ inv_sig
        except Exception:
            Cf = apply_rmt_filter(r)
            Sf = np.outer(sig, sig) * Cf

        adj     = (np.abs(Cf) > 0.5).astype(int)
        G       = nx.from_numpy_array(adj)
        density = nx.density(G)

        # Gamma adaptif — dibatasi gamma_max untuk profil defensif
        gamma_t = np.clip(density * self.sensitivity, 0.1, self.gamma_max)

        dist  = build_mst(Cf)
        cent  = compute_eigenvector_centrality(dist)
        w_net = 1.0 / (cent + 1e-6)
        w_net /= w_net.sum()

        res_mko = minimize(
            lambda w: w @ Sf @ w,
            np.ones(n) / n,
            method='SLSQP',
            bounds=[(0, 1)] * n,
            constraints=({'type': 'eq', 'fun': lambda x: np.sum(x) - 1})
        )
        return (1 - gamma_t) * res_mko.x + gamma_t * w_net

# 6. ML-Gated Dynamic Gamma (Model Usulan Lanjutan)
class MLGatedNetworkStrategy(PortfolioStrategy):
    """
    Strategi NW dengan Dynamic Gamma. Menggunakan classifier ML (SVM/LightGBM)
    untuk mendeteksi market regime secara otomatis.
    """
    def __init__(self, name='ML-Gated NW', model=None):
        super().__init__(name)
        self.model = model # Pre-trained classifier (e.g., SVM/LightGBM)

    def get_weights(self, r):
        # 1. Feature Engineering (Volatility, Momentum, etc.)
        features = np.array([r.mean().mean(), r.std().mean()]).reshape(1, -1)
        
        # 2. Predict Market Phase (0: Bearish, 1: Bullish, 2: Sideways)
        # Jika model belum dilatih, default ke NW standar
        phase = self.model.predict(features)[0] if self.model else 0
        
        # 3. Dynamic Gamma Setting
        # Bearish: Gamma tinggi (fokus ketahanan jaringan)
        # Bullish: Gamma rendah (fokus ke efisiensi Markowitz)
        gamma = 1.0 if phase == 0 else (0.1 if phase == 1 else 0.5)
        
        # 4. NW Optimization with dynamic gamma
        # (Logika serupa dengan NetworkMarkowitz namun dengan gamma dari ML)
        return NetworkMarkowitz(gamma=gamma).get_weights(r)
\end{lstlisting}
\textbf{Logika Implementasi:}
\begin{itemize}
    \item \textbf{Base Class:} Mendefinisikan \textit{blueprint} \texttt{get\_weights(returns\_data)} yang harus diimplementasikan oleh setiap subclass.
    \item \textbf{Mean-Variance Optimization:} Menggunakan \texttt{scipy.optimize.minimize} dengan fungsi objektif kuadratik $w^T \Sigma w$ dan kendala \texttt{equality} (jumlah bobot = 1) serta \texttt{bounds} (0, 1) untuk strategi \textit{long-only}.
    \item \textbf{Network Penalty:} Pada strategi NW, fungsi objektif dimodifikasi menjadi $f(w) = w^T \Sigma w + \gamma \sum (\text{centrality}_i \cdot w_i)$.
    \item \textbf{AGGP Dynamic Gate:} Menghitung densitas jaringan $D$. Bobot akhir dihitung sebagai kombinasi linear: $w_{\text{final}} = (1 - \gamma_t) w_{\text{mko}} + \gamma_t w_{\text{net}}$, di mana $\gamma_t$ di-\textit{clip} secara dinamis.
    \item \textbf{ML-Gated Dynamic Gamma:} Menggunakan model klasifikasi (SVM/LightGBM) untuk mendeteksi rezim pasar (\textit{Bullish, Bearish, Sideways}). Parameter $\gamma$ diatur secara otomatis ($\gamma=1.0$ saat krisis, $\gamma=0.1$ saat stabil) untuk menyeimbangkan diversifikasi jaringan dan efisiensi rata-rata.
\end{itemize}

\section{Simulasi dan Evaluasi (Backtesting)}

\subsection{Sel 8 — Framework Backtesting}
\textbf{Deskripsi:} Pengembangan mesin simulasi perdagangan (\textit{trading engine}) untuk menguji keandalan strategi pada data historis.
\textbf{Isi Kode:}
\begin{lstlisting}
def backtest_strategy(strategy, df_returns, window_size=120, rebalance_freq=7, transaction_cost=0.001):
    """Simulasi backtest dengan rolling window dan biaya transaksi."""
    portfolio_returns = []
    dates             = []

    for i in range(window_size, len(df_returns), rebalance_freq):
        train = df_returns.iloc[i - window_size:i]
        w     = strategy.get_weights(train)

        test_end  = min(i + rebalance_freq, len(df_returns))
        test_data = df_returns.iloc[i:test_end]

        for j in range(len(test_data)):
            daily_ret = np.dot(w, test_data.iloc[j].values)
            if j == 0 and len(portfolio_returns) > 0:
                daily_ret -= transaction_cost
            portfolio_returns.append(daily_ret)
            dates.append(test_data.index[j])

    res_df = pd.DataFrame({'date': dates, 'return': portfolio_returns})
    res_df['cumulative_return'] = (1 + res_df['return']).cumprod()

    return {
        'strategy':           strategy.name,
        'returns':            np.array(portfolio_returns),
        'cumulative_returns': res_df['cumulative_return'].values,
        'results_df':         res_df
    }
\end{lstlisting}
\textbf{Logika Kode:}
\begin{itemize}
    \item \textbf{Rolling Loop:} Fungsi \texttt{backtest\_strategy} menggunakan \textit{loop} yang bergeser sejauh \textit{rebalance\_freq}.
    \item \textbf{Out-of-Sample:} Bobot dihitung pada jendela \textit{train} (120 hari) dan diaplikasikan pada jendela \textit{test} berikutnya (7 hari).
    \item \textbf{Transaction Costs:} Mengurangi return harian sebesar 0.1\% dikalikan pergantian (\textit{turnover}) pada setiap hari rebalancing.
    \item \textbf{Compounding:} Return kumulatif dihitung menggunakan \texttt{.cumprod()} untuk mensimulasikan efek bunga majemuk.
\end{itemize}

\subsection{Sel 9 — Eksekusi Backtesting}
\textbf{Deskripsi:} Tahap komputasi untuk menjalankan seluruh strategi secara paralel.
\textbf{Isi Kode:}
\begin{lstlisting}
# --- Inisialisasi semua strategi ---
gamma_values = [0.005, 0.025, 0.05, 0.15, 0.7, 1.0]

strategies = (
    [
        EquallyWeighted("EW"),
        ClassicalMarkowitz("CM"),
        GlassoMarkowitz("GM"),
        NetworkMarkowitz("NW (gamma=0)", gamma=0),
    ]
    + [NetworkMarkowitz(f"NW (gamma={g})", gamma=g) for g in gamma_values]
    + [
        GraphDiversification('Graph Divers. (theta=0.4)', corr_threshold=0.4),
        GraphDiversification('Graph Divers. (theta=0.5)', corr_threshold=0.5),
        RMTGraphDiversification('RMT GD (theta=0.4)',     corr_threshold=0.4),
        AdaptiveGraphPortfolio('AGGP (Optimized)',         sensitivity=2.0),
    ]
)

# --- Eksekusi backtest ---
results = {}
for strat in strategies:
    print(f"Running: {strat.name} ...")
    results[strat.name] = backtest_strategy(strat, df_returns)

print("\nAll backtests completed!")
\end{lstlisting}
\textbf{Operasional:} Inisialisasi daftar objek strategi (termasuk variasi $\gamma$) dan penyimpanan hasil simulasi ke dalam \textit{dictionary} \texttt{results} untuk analisis statistik lebih lanjut.

\section{Hasil dan Pembahasan (Results)}

\subsection{Sel 10 — Analisis Performa Periodik (Table 2)}
\textbf{Isi Kode:}
\begin{lstlisting}
# Identifikasi tanggal pelaporan (Januari, Mei, September)
report_dates = ['2018-01-31', '2018-05-31', '2018-09-30', '2019-01-31', '2019-05-31', '2019-09-30']

table_data = []
for name, res in results.items():
    row = {'Strategy': name}
    df = res['results_df'].set_index('date')
    for date in report_dates:
        try:
            val = df.iloc[df.index.get_indexer([date], method='pad')[0]]['cumulative_return']
            row[date] = f"{(val-1)*100:.2f}%"
        except:
            row[date] = "N/A"
    table_data.append(row)

print(pd.DataFrame(table_data).set_index('Strategy'))
\end{lstlisting}

\subsection{Sel 11 — Visualisasi Performa Kumulatif (Figure 6)}
\textbf{Isi Kode:}
\begin{lstlisting}
plt.figure(figsize=(12, 7))
for name, res in results.items():
    if name in ['EW', 'CM', 'GM', 'NW (gamma=1.0)', 'AGGP (Optimized)']:
        plt.plot(res['results_df']['date'], res['results_df']['cumulative_return']*100, label=name)

plt.title('Figure 6 | Cumulative Portfolio Performance ($100 Investment)')
plt.ylabel('Portfolio Value ($)')
plt.legend()
plt.show()
\end{lstlisting}
\textbf{Logika Kode:} Plotting dilakukan terhadap \textit{baseline} \$100. Fungsi ini mengompilasi semua seri data dari \texttt{results\_df} ke dalam koordinat sumbu-y yang seragam guna mendeteksi koherensi dan divergensi performa antar strategi.

\subsection{Sel 12 \& 13 — Metrik Performa Risiko (Table 3 \& 4)}
\textbf{Isi Kode:}
\begin{lstlisting}
# Sharpe and VaR Analysis Loop
risk_metrics = []
for name, res in results.items():
    rets = res['returns']
    ann_return = np.mean(rets) * 365
    ann_vol    = np.std(rets) * np.sqrt(365)
    sharpe     = ann_return / (ann_vol + 1e-8)
    var95      = calculate_var(rets, 0.95)
    
    risk_metrics.append({
        'Strategy': name, 'Sharpe': sharpe, 'VaR 95%': var95
    })
print(pd.DataFrame(risk_metrics).set_index('Strategy'))
\end{lstlisting}
\textbf{Logika Kode:}
\begin{itemize}
    \item \textbf{Rolling Sharpe:} Menggunakan \texttt{pd.DateOffset(months=4)} untuk menghitung volatilitas dan pengembalian rata-rata secara dinamis.
    \item \textbf{VaR (95\%):} Dihitung menggunakan persentil ke-5 dari distribusi log-return harian untuk memodelkan skenario terburuk harian.
\end{itemize}

\subsection{Sel 14 — Analisis Tail-Risk: Rachev Ratio (Table 5)}
\textbf{Isi Kode:}
\begin{lstlisting}
table5_data = []
for name, res in results.items():
    rr = calculate_rachev_ratio(res['returns'])
    table5_data.append({'Strategy': name, 'Rachev': round(rr, 4)})
print(pd.DataFrame(table5_data))
\end{lstlisting}
\textbf{Logika Kode:} Fokus pada integrasi \texttt{calculate\_rachev\_ratio} dengan parameter $\alpha=10\%$. Logika ini sangat krusial untuk membedakan antara model Markowitz standar dengan model berbasis jaringan (NW/AGGP) dalam menangani distribusi \textit{fat-tail}.
        
\subsection{Sel 15 — Analisis Fase Pasar (Table 6)}
\textbf{Isi Kode:}
\begin{lstlisting}
# Definisi fase pasar
fase_pasar = {
    'Bearish':   ('2018-01-01', '2019-03-31'),
    'Recovery':  ('2019-04-01', '2019-06-30'),
    'Stable':    ('2019-07-01', '2019-10-17')
}

phase_strategies = ['EW', 'CM', 'GM', 'NW (gamma=0)', 'NW (gamma=0.15)', 'AGGP (Optimized)']

# --- Panel A: Sharpe Ratio per Fase ---
rows_sr = []
for f_name, (start, end) in fase_pasar.items():
    row = {'Market Phase': f_name}
    for s_name in phase_strategies:
        df = results[s_name]['results_df']
        mask = (df['date'] >= start) & (df['date'] <= end)
        rets = df.loc[mask, 'return']
        sr = (np.mean(rets) / np.std(rets)) * np.sqrt(365) if not rets.empty else 0
        row[s_name] = sr
    rows_sr.append(row)

print("TABLE 6 | Panel A: Sharpe Ratio per Market Phase")
print(pd.DataFrame(rows_sr).set_index('Market Phase'))

# --- Panel B: Rachev Ratio per Fase ---
rows_rr = []
for f_name, (start, end) in fase_pasar.items():
    row = {'Market Phase': f_name}
    for s_name in phase_strategies:
        df = results[s_name]['results_df']
        mask = (df['date'] >= start) & (df['date'] <= end)
        rets = df.loc[mask, 'return'].values
        rr = calculate_rachev_ratio(rets) if len(rets) > 0 else 0
        row[s_name] = rr
    rows_rr.append(row)

print("\nTABLE 6 | Panel B: Rachev Ratio per Market Phase")
print(pd.DataFrame(rows_rr).set_index('Market Phase'))
\end{lstlisting}
\textbf{Logika Kode:} Melakukan segmentasi DataFrame berdasarkan \textit{dictionary} \texttt{fase\_pasar}. Statistik Sharpe dan Rachev dihitung ulang hanya untuk data dalam rentang tanggal fase \textit{Bearish}, \textit{Recovery}, dan \textit{Stable} guna membuktikan kemampuan adaptasi model AGGP.

\section{Kesimpulan}
Penelitian ini berhasil membuktikan bahwa integrasi metrik jaringan (\textit{Network Analysis}) ke dalam optimasi portofolio kripto memberikan ketahanan yang signifikan terhadap risiko sistemik. Model usulan \textbf{Adaptive Graph-Gated Portfolio (AGGP)} menunjukkan performa paling optimal sebagai model "segala medan" yang mampu menyesuaikan postur investasi secara otomatis berdasarkan dinamika korelasi pasar.

\end{document}
